\documentclass[11pt]{article}
\usepackage{amsmath,amssymb,amsthm}
\usepackage[margin=1.1in]{geometry}
\newtheorem{lemma}{Lemma}
\newtheorem{proposition}{Proposition}
\newtheorem{theorem}{Theorem}
\theoremstyle{remark}
\newtheorem{remark}{Remark}
\newcommand{\li}{\operatorname{li}}
\newcommand{\Z}{\mathbb{Z}}

\title{Certified exact sums of primes over dyadic intervals\\ via a ladder of generalized factorials}
\author{Carl Gribble\thanks{Draft v0.1, companion to \emph{Prime counts and prime sums from multiplication tables: a
self-similar formula system}~\cite{Companion} (same author). Developed in collaboration with Claude
(Anthropic); all numerical claims were machine-verified as described in Section~5. Responsibility for the content rests with the author.
A full disclosure of AI use appears at the end of the paper.}}
\date{August 2026 --- draft}

\begin{document}
\maketitle

\begin{abstract}
For an integer $m \ge 8$ let $S(m) = \sum_{m < p \le 2m} p$, the sum over primes in the dyadic window
$(m, 2m]$. We construct a ladder of generalized factorials
$A_j(2m) = \prod_{i \le 2m} i^{\,Q_j(i)}$ whose integer exponent sequences
$Q_j(i) = m^j\,T_j\!\big(\tfrac{2i-3m}{m}\big)$ are Chebyshev polynomials rescaled to take integer values.
The $p$-adic valuation of $A_j(2m)$ at a prime $p \in (m,2m]$ equals $Q_j(p)$, so each rung of the ladder
yields one exact linear functional (``moment'') of the interval's primes, computable from the primes
$\le m$ together with logarithms of the integers $\le 2m$, without identifying any prime inside the
interval. A near-minimax Chebyshev combination of $K{+}1$ moments approximates $t/\log t$; a rigorous
sandwich, weighted by the zeroth moment, encloses $S(m)$ in an interval of certified width. Since $S(m)$
is an integer, a certified width below $1$ determines it exactly. An implementation in $260$-bit interval
arithmetic certifies $S(10^3) = 200{,}923$, $S(10^4) = 15{,}434{,}795$ and
$S(10^5) = 1{,}255{,}204{,}276$ using $6$, $8$ and $10$ moments respectively, with empirical window decay
of roughly one order of magnitude per additional moment, above the proven geometric floor $3 + 2\sqrt{2}$.
We situate the construction relative to Chebyshev-type identities and to combinatorial prime-summing
algorithms, and we state its limitations plainly: the method is a certified exact procedure of
logarithmic description length, not an asymptotically fast algorithm and not a closed form.
\end{abstract}

\medskip
\noindent\textbf{Keywords:} primes in dyadic intervals; Chebyshev's method; generalized factorials;
$p$-adic valuations; Chebyshev polynomials; certified computation; interval arithmetic.

\smallskip
\noindent\textbf{MSC 2020:} primary 11N05; secondary 11Y35, 41A50, 65G30.

\section{Introduction}

Exact computation of prime-counting and prime-summing functions has two established traditions. The
combinatorial tradition descends from Meissel through Lehmer, Lagarias--Miller--Odlyzko and
Del\'eglise--Rivat, and computes $\pi(x)$ in $O(x^{2/3}\log^{-2}x)$ time \cite{DR,Staple}; it has been
extended to the prime power sums $\sum_{p \le x} p^k$ \cite{Orlov}, with widely used implementations
\cite{primesum}, and has a folklore branch, the $O(x^{3/4})$ dynamic program posted by the Project Euler
user Lucy\_Hedgehog in 2013, which sums the primes below $x$ ``without finding them all''
\cite{Lucy}. The analytic tradition, initiated by Lagarias and Odlyzko and realized rigorously by Platt
\cite{Platt} and B\"uthe \cite{Buthe}, evaluates an explicit formula over zeta zeros to accuracy below
$\tfrac12$ and rounds; the target being an integer converts a certified approximation into an exact
determination.

This note describes a third route to exact interval prime sums, elementary in its ingredients and
certified in its execution. It originates in Chebyshev's 1852 observation \cite{Cheb} that the prime
factorization of the central binomial coefficient controls $\theta(2m)-\theta(m)$: the primes of
$(m, 2m]$ appear in $\binom{2m}{m}$ with exponent exactly one. We extend this single exact functional to
a \emph{ladder} of them. The factorial is the weight-$1$ member of the family
$\prod_{i \le n} i^{w(i)}$; taking polynomial weights $w$ produces generalized factorials (the weight
$w(i)=i$ gives the classical hyperfactorial \cite{HypVal1,HypVal2}) whose valuations at large primes are
$w(p)$ itself. Choosing the weights to be integer-rescaled Chebyshev polynomials yields a numerically
stable basis of exact moments; approximation theory then converts finitely many moments into a rigorous
two-sided enclosure of $S(m)$, and the integrality of $S(m)$ converts a sufficiently narrow enclosure
into the exact value.

We make three claims and two explicit non-claims. Claimed: (i) the exact extraction identities of
Section~2; (ii) the certified enclosure protocol of Sections~3--4, whose window width contracts at a
guaranteed geometric rate in the number of moments; (iii) certified computations at three scales,
verified against sieves. Not claimed: computational superiority over the combinatorial algorithms, which
remain asymptotically faster; and novelty of the individual ingredients, each of which is classical. To
the best of our knowledge the assembly --- polynomial-weight generalized factorials organized as a moment
ladder and closed into a certified integer determination --- has not appeared in the literature, but we
regard this as a statement about our searches, not a theorem.

\section{The ladder and its exact moments}

Fix an integer $m \ge 8$ and write $n = 2m$, $u(t) = (2t - 3m)/m$, so that $u$ maps $[m, 2m]$ onto
$[-1, 1]$. Let $T_j$ denote the Chebyshev polynomials of the first kind and define integer sequences
\[
Q_j(i) \;=\; m^{j}\, T_j\big(u(i)\big), \qquad j \ge 0 .
\]
That $Q_j(i) \in \Z$ for $i \in \Z$ follows from the integer recurrence
$Q_0 = 1$, $Q_1(i) = 2i - 3m$, $Q_{j+1}(i) = 2(2i-3m)\,Q_j(i) - m^2 Q_{j-1}(i)$.
Define the $j$-th ladder object
\[
A_j(n) \;=\; \prod_{i=1}^{n} i^{\,Q_j(i)} \in \mathbb{Q}_{>0},
\]
a rational number in general since $Q_j$ takes both signs; $A_0(n) = n!$ and, for the weight $i$ in
place of $Q_1$, one obtains the hyperfactorial. The ladder interacts with primes through the following
elementary facts.

\begin{lemma}[Interval valuations]\label{lem:val}
Let $p$ be prime with $m < p \le 2m$. Then $v_p\big(A_j(n)\big) = Q_j(p)$.
\end{lemma}

\begin{proof}
Since $p > m = n/2$, the only multiple of $p$ in $[1, n]$ is $p$ itself, and $p^2 > n$ for $m \ge 2$.
Hence $v_p(A_j(n)) = \sum_{i \le n} Q_j(i)\, v_p(i) = Q_j(p)$.
\end{proof}

\begin{lemma}[Below-interval valuations]\label{lem:corr}
Let $p \le m$ be prime. Then
\[
v_p\big(A_j(n)\big) \;=\; \sum_{a \ge 1,\; p^a \le n} \;\sum_{t=1}^{\lfloor n/p^a \rfloor} Q_j\big(p^a t\big),
\]
and each inner sum is an exact integer computable in $O(j^2)$ arithmetic operations from the coefficient
list of $Q_j$, via the exact power sums
$\sum_{t=1}^{T} t^r = \sum_{q=1}^{r} q!\, S(r,q) \binom{T+1}{q+1}$,
where $S(r,q)$ are Stirling numbers of the second kind.
\end{lemma}

\begin{proof}
The double sum is $\sum_{i \le n} Q_j(i) v_p(i)$ rewritten by counting, for each prime power $p^a$, the
indices $i$ it divides. Substituting $i = p^a t$ makes the inner sum a polynomial in $t$ summed over an
initial segment, which the stated Faulhaber-type identity evaluates exactly.
\end{proof}

\begin{proposition}[Exact moment extraction]\label{prop:extract}
For every $j \ge 0$,
\[
M_j \;:=\; \sum_{m < p \le 2m} Q_j(p) \log p
\;=\; \sum_{i=2}^{n} Q_j(i) \log i \;-\; \sum_{p \le m} v_p\big(A_j(n)\big)\, \log p ,
\]
with the valuations given exactly by Lemma~\ref{lem:corr}. The right-hand side involves only the primes
$\le m$ and logarithms of integers $\le n$; no integer in $(m, 2m]$ is examined for primality.
\end{proposition}

\begin{proof}
Take logarithms in the prime factorization of $A_j(n)$, split the primes at $m$, and apply
Lemma~\ref{lem:val} to the upper range. Prime powers $p^a \le n$ with $a \ge 2$ have $p \le \sqrt{n} \le m$
and are absorbed by Lemma~\ref{lem:corr}.
\end{proof}

The case $j = 0$ is Chebyshev's identity in the form
$\theta(2m) - \theta(m) = \log n! - \sum_{p \le m} v_p(n!) \log p$; we write $\mu_0 = M_0$.

\section{From moments to the integer}

Let $f(t) = t/\log t$, so that $f(p)\log p = p$ and
$S(m) = \sum_{m < p \le 2m} f(p) \log p$. Fix coefficients $a_0, \dots, a_K \in \mathbb{R}$
(in practice, dyadic rationals), set $P(t) = \sum_{j \le K} a_j T_j(u(t))$ and $g = f - P$ on
$[m, 2m]$. Since $Q_j(p) = m^j T_j(u(p))$,
\begin{equation}\label{eq:sandwich}
S(m) \;=\; \sum_{j=0}^{K} a_j\, m^{-j} M_j \;+\; \sum_{m < p \le 2m} g(p) \log p ,
\qquad
\mu_0 \min_{[m,2m]} g \;\le\; \sum_p g(p)\log p \;\le\; \mu_0 \max_{[m,2m]} g ,
\end{equation}
the bounds valid because every $\log p > 0$ and $\sum_p \log p = \mu_0$. Thus any choice of $P$ yields a
rigorous enclosure of $S(m)$ of width $\mu_0 \cdot \mathrm{osc}_{[m,2m]}(g)$; when the certified width is
below $1$, the unique integer in the enclosure is $S(m)$. Validity does not depend on optimality of $P$,
which may therefore be designed in floating point and then frozen.

Two elementary bounds certify the range of $g$ from finitely many evaluations.

\begin{lemma}[Curvature bounds]\label{lem:curv}
For $t \ge e^2$ one has $|f''(t)| = \big|\tfrac{2}{\log t} - 1\big| \big/ (t \log^2 t) \le
1/(m \log^2 m)$ on $[m, 2m]$, and $\max_{[-1,1]} |T_j''| = j^2(j^2-1)/3$. Consequently
$\|g''\|_\infty \le B_2 := \frac{1}{m \log^2 m} + \frac{4}{m^2} \sum_{j \le K} |a_j|\,
\frac{j^2(j^2-1)}{3}$, and for a uniform grid of step $s$ containing both endpoints,
\[
\max_{[m,2m]} g \;\le\; \max_i g(t_i) + \tfrac{s^2}{8} B_2 , \qquad
\min_{[m,2m]} g \;\ge\; \min_i g(t_i) - \tfrac{s^2}{8} B_2 .
\]
\end{lemma}

\begin{proposition}[Geometric window decay]\label{prop:rate}
Let $E_K$ denote the error of best uniform approximation to $f$ on $[m,2m]$ by polynomials of degree
$K$. Then $\limsup_{K \to \infty} E_K^{1/K} \le 1/\rho_m$, where $\rho_m$ is the Bernstein-ellipse
parameter of the nearest singularity of $f$ ($t = 1$, where $\log t = 0$) relative to $[m,2m]$; as
$m \to \infty$, $\rho_m \to 3 + 2\sqrt{2} \approx 5.83$. Hence the achievable certified window
contracts at least geometrically in $K$, and the number of moments needed to reach width $1$ grows only
logarithmically in the initial window.
\end{proposition}

\begin{proof}[Proof sketch]
$f$ is analytic in a neighborhood of $[m, 2m]$ with nearest singularities at $t = 1$ and $t = 0$; under
$u(t)$ these map near $u = -3$, and $\rho = |u_\ast| + \sqrt{u_\ast^2 - 1}$ with $u_\ast \to -3$. The
bound is then the classical Bernstein estimate; see \cite[Ch.~8]{Tref}.
\end{proof}

\begin{remark}
Empirically the contraction factor observed in Section~5 is close to $10$ per moment, faster than the
guaranteed floor; the floor alone already gives the logarithmic moment count.
\end{remark}

\section{Certified implementation}

The protocol has four stages (reference implementation in the code repository \cite{Code}), all
executed in ball (interval) arithmetic at $260$ bits except where exact integer arithmetic is
available. (1) \emph{Moments.} The anchor $\sum_{i \le n} Q_j(i) \log i$ is
accumulated with $Q_j(i)$ produced by the integer recurrence and $\log i$ as an interval; the
corrections of Lemma~\ref{lem:corr} are exact integers multiplied by interval logarithms. This yields
enclosures of $M_0, \dots, M_K$ whose observed widths at the scales below were at most
$4.2 \times 10^{-11}$ (and as small as $10^{-71}$). (2) \emph{Design.} $a_0, \dots, a_K$ are Chebyshev
interpolation coefficients of $f$ computed in floating point and frozen exactly as dyadics. (3)
\emph{Residual certification.} $g$ is evaluated in interval arithmetic on an integer grid of step $s$
and extended to $[m, 2m]$ by Lemma~\ref{lem:curv}. (4) \emph{Assembly.} Equation~\eqref{eq:sandwich} is
evaluated in interval arithmetic; if $\lceil \text{lo} \rceil = \lfloor \text{hi} \rfloor$, that integer
is $S(m)$.

The trust base is: correctness of the interval arithmetic library (mpmath), of big-integer arithmetic,
and of Lemma~\ref{lem:curv}'s two analytic inequalities. All other quantities are enclosed, not
estimated. Ground-truth verification by segmented sieve, and cross-checks of every moment enclosure
against direct computation over the (sieved) interval primes, are performed outside the pipeline.

\section{Results}

\begin{center}
\begin{tabular}{lcccc}
\hline
window & moments $K{+}1$ & certified width & $S(m)$ determined & sieve check \\
\hline
$(10^3,\,2\cdot 10^3]$ & $6$ & $0.140$ & $200{,}923$ & agrees \\
$(10^4,\,2\cdot 10^4]$ & $8$ & $0.155$ & $15{,}434{,}795$ & agrees \\
$(10^5,\,2\cdot 10^5]$ & $10$ & $0.253$ & $1{,}255{,}204{,}276$ & agrees \\
\hline
\end{tabular}
\end{center}

\noindent
Runtimes on commodity hardware were $0.3$, $3.1$ and $31$ seconds respectively, dominated by the
$O(m)$ interval logarithms of the anchors. Measured window contraction per added moment was $9.7$--$11.4$
across the three scales. In uncertified experiments the same estimator tracked sieved interval sums with
relative errors of order $10^{-4}$ at additional scales.

\section{Related identities and remarks}

Rung $0$ is Chebyshev's binomial identity \cite{Cheb}. The weight-$i$ rung is governed by the
hyperfactorial, whose $p$-adic valuations have an independent literature \cite{HypVal1,HypVal2}; the
exponent formula involves triangular numbers, and more generally the corrections of
Lemma~\ref{lem:corr} are the Faulhaber closed forms of iterated-sum sequences, so the ladder may be
read as a multiplicative avatar of the tower $1, n, \binom{n+1}{2}, \binom{n+2}{3}, \dots$. High-precision
evaluation of the anchors by Euler--Maclaurin summation leads to the generalized Glaisher--Kinkelin
(Bendersky--Adamchik) constants, expressible through $\zeta'(-k)$ \cite{Bend,Adam,Coppo}; we have not
needed this refinement at the scales reported, but it removes the $O(m)$ anchor cost in principle. A
prime-free variant of the zeroth-order theory is provided by Linnik's identity \cite{Linnik}, which
expresses $\Lambda(n)/\log n$ as an alternating-harmonic combination of the counts of ordered
factorizations into parts $\ge 2$, and was developed into prime-counting and prime-power-sum
algorithms by McKenzie \cite{McKenzie}; the companion paper \cite{Companion} builds this route into a
uniform weighted formula system which, in exact rational arithmetic, reproduces
$\sum_{p \le 10^5} p = 454{,}396{,}537$ from divisor-lattice data alone. Finally, the dyadic windows
$(2^{r-1}, 2^r]$ tile the integers and each contains a prime by Bertrand's postulate, so the pair
$\big(I(r), F\big)$ with $I(r) = (2^{r-1}, 2^r]$ and $F$ the certified procedure above answers, for
every $r$, the question ``what is the sum of the primes in the $r$-th dyadic window'' exactly.

\section{Limitations and open questions}

The method as implemented costs $O(m)$ interval logarithms and is therefore not competitive with the
$O(x^{2/3})$ combinatorial algorithms \cite{DR,Orlov} as a computational tool; its interest is
methodological. The corrections consume the primes below the window; a fully self-contained variant
recurses dyadically to a trivial base, or replaces rung $0$ by the Linnik route
\cite{McKenzie,Companion}. Three questions seem
worth recording. (1) Can the anchors be evaluated to certified accuracy in time polylogarithmic in $m$
via the Euler--Maclaurin/$\zeta'(-k)$ route, making the whole pipeline sublinear? (2) Does the ladder
yield explicit \emph{inequalities} for $S(m)$ of Chebyshev strength but higher order --- bounds rather
than computations --- that the combinatorial algorithms do not provide? (3) Is there a principled
obstruction theorem formalizing the folklore that no fixed-length elementary closed form for $S(m)$
exists, given that the fluctuating part of $S(m)$ carries the zeta spectrum? We do not know, and we
would welcome pointers to prior appearances of this construction.

\section*{Disclosure of AI use}

In accordance with publisher policies on the use of artificial intelligence, the author discloses the
following. This paper was developed in an extensive collaboration with Claude (Anthropic; Claude Fable
5, accessed through the claude.ai interface, 2026). The direction of the work --- the founding research
programme of exact moment systems for interval primes, of which this paper is the second --- together
with all decisions about scope, claims, and submission, are the author's. Within that direction, the
AI system contributed substantially to the mathematical development and proofs, wrote the verification
code, drafted and revised the manuscript text, and assisted with literature search; the reason for its
use was to enable a single non-specialist author to develop, check, and document the material to a
publishable standard. Every numerical and mathematical claim was subsequently machine-verified by the
certified computations reported in Section~5, with the verification code openly available in the cited
repository \cite{Code}. The author has reviewed the entire manuscript and takes full responsibility
for its content, including the accuracy of the mathematics and of all references.

\begin{thebibliography}{99}
\bibitem{Cheb} P.~L. Chebyshev, M\'emoire sur les nombres premiers, \emph{J. Math. Pures Appl.}
\textbf{17} (1852), 366--390.
\bibitem{DR} M. Del\'eglise and J. Rivat, Computing $\pi(x)$: the Meissel, Lehmer, Lagarias, Miller,
Odlyzko method, \emph{Math. Comp.} \textbf{65} (1996), 235--245.
\bibitem{Staple} D.~B. Staple, The combinatorial algorithm for computing $\pi(x)$, arXiv:1503.01839
(2015).
\bibitem{Orlov} A. Orlov, The combinatorial method to compute the sum of the powers of primes,
arXiv:2111.15545 (2021).
\bibitem{primesum} K. Walisch, \emph{primesum} (software), \texttt{github.com/kimwalisch/primesum}.
\bibitem{Lucy} Lucy\_Hedgehog, post in the Project Euler Problem 10 discussion thread (2013); see also
G. Broxey, ``Lucy's Algorithm + Fenwick Trees'' (2023),
\texttt{https://gbroxey.github.io/blog/2023/04/09/lucy-fenwick.html}.
\bibitem{Platt} D.~J. Platt, Computing $\pi(x)$ analytically, \emph{Math. Comp.} \textbf{84} (2015),
1521--1535.
\bibitem{Buthe} J. B\"uthe, An analytic method for bounding $\psi(x)$, \emph{Math. Comp.} \textbf{87}
(2018); An improved analytic method for calculating $\pi(x)$, \emph{Manuscripta Math.} \textbf{151}
(2016), 329--352.
\bibitem{HypVal1} L. Onnis, On the $p$-adic valuation of a hyperfactorial, arXiv:2109.05616 (2021).
\bibitem{HypVal2} J.-C. Pain, Bounds on the $p$-adic valuation of the factorial, hyperfactorial and
superfactorial, arXiv:2408.00353 (2024).
\bibitem{Bend} L. Bendersky, Sur la fonction gamma g\'en\'eralis\'ee, \emph{Acta Math.} \textbf{61}
(1933), 263--322.
\bibitem{Adam} V.~S. Adamchik, Polygamma functions of negative order, \emph{J. Comput. Appl. Math.}
\textbf{100} (1998), 191--199.
\bibitem{Coppo} M.-A. Coppo, Generalized Glaisher--Kinkelin constants and Ramanujan summation of
series, \emph{Res. Number Theory} \textbf{10} (2024), art.~8.
\bibitem{Linnik} Yu.~V. Linnik, \emph{The Dispersion Method in Binary Additive Problems}, Amer. Math.
Soc., Providence, 1963.
\bibitem{Tref} L.~N. Trefethen, \emph{Approximation Theory and Approximation Practice}, SIAM,
Philadelphia, 2013.
\bibitem{McKenzie} N. McKenzie, Computing the prime counting function with Linnik's identity,
manuscript (2011), available at \texttt{https://www.primecounting.com}.
\bibitem{Companion} C. Gribble, Prime counts and prime sums from multiplication tables: a
self-similar formula system, companion manuscript, Zenodo (2026), \texttt{doi:10.5281/zenodo.21769103}.
\bibitem{Code} C. Gribble, Reference implementations accompanying this paper and \cite{Companion},
code repository (2026), \texttt{https://github.com/carlgribble-caa/prime-moments}.
\end{thebibliography}

\end{document}
